\chapter{INTRODUÇÃO}
\label{cap:introducao}

\section{Contextualização e Cenário da Pesquisa}
\label{sec:contextualizacao}

A segurança do paciente consolidou-se, nas últimas décadas, como um dos pilares mais urgentes e desafiadores da saúde pública global. Dados consolidados pela Organização Mundial da Saúde (OMS) apontam que, em escala mundial, um em cada dez pacientes hospitalizados sofre algum tipo de dano prevenível decorrente da assistência à saúde, resultando em milhões de óbitos anuais e em custos econômicos que superam centenas de bilhões de dólares \cite{who2019patient, who2021global}. No Brasil, o Ministério da Saúde e a Agência Nacional de Vigilância Sanitária (ANVISA) reforçam que os incidentes assistenciais e os Eventos Adversos (EAs) constituem expressiva parcela das causas de prolongamento de internação hospitalar, sequelas temporárias ou permanentes e ocupação indevida de leitos críticos \cite{brasil2013pnsp}.

Historicamente, o monitoramento e a vigilância de incidentes em serviços de saúde têm sido fundamentados em sistemas de notificação espontânea e voluntária por parte da equipe assistencial. Não obstante sua importância institucional, a literatura médica e de enfermagem demonstra exaustivamente que a notificação voluntária padece de grave subnotificação, capturando comprovadamente menos de 10\% dos eventos adversos que efetivamente atingem os pacientes \cite{griffin2009ihi, classen2011adverse}. Fatores como a cultura punitiva persistente, a sobrecarga de trabalho dos profissionais na beira do leito, a ausência de retorno formativo (\textit{feedback}) e a complexidade burocrática dos formulários tradicionais desestimulam o registro tempestivo dos incidentes.

Diante desse cenário e da necessidade premente de metodologias objetivas e reprodutíveis, o \textit{Institute for Healthcare Improvement} (IHI), nos Estados Unidos, formulou a metodologia \textit{Global Trigger Tool} (IHI-GTT) \cite{griffin2009ihi}. Diferentemente dos sistemas reativos e passivos, a GTT preconiza uma abordagem ativa, sistemática e retrospectiva de revisão de prontuários com base em ``rastreadores'' ou ``gatilhos'' clínicos padronizados. Esses gatilhos (como o uso inadvertido de antídotos específicos como a naloxona, quedas abruptas nos níveis séricos de hemoglobina, extubações acidentais na Unidade de Terapia Intensiva ou transferências não planejadas para setores de maior complexidade) funcionam como sentinelas, orientando a equipe de auditores diretamente aos episódios com elevada probabilidade de ocorrência de dano assistencial.

No âmbito da Universidade Federal de Sergipe (UFS), a inserção do Hospital Universitário (HU-UFS/EBSERH) como centro de referência de média e alta complexidade, aliado aos cursos de graduação e pós-graduação das áreas de Ciências Exatas e Saúde, cria um ecossistema propício para a inovação interdisciplinar. Todavia, a operacionalização prática da metodologia GTT nas instituições de ensino e saúde brasileiras ainda é realizada de forma predominantemente artesanal, com planilhas eletrônicas desconexas, formulários em papel e ausência de plataformas educacionais dinâmicas que permitam aos futuros profissionais de saúde aprenderem e aplicarem a metodologia de forma prática e supervisionada.

\section{Finalidade do Projeto}
\label{sec:finalidade}

A finalidade primordial deste projeto compreende a construção de um sistema de monitoramento de eventos adversos como ferramenta educacional no ensino de gerenciamento de riscos e segurança do paciente hospitalizado, estabelecendo um suporte interativo e eficaz no processo de ensino-aprendizagem e na formação de profissionais de enfermagem, com vistas a integrar a segurança do paciente com a educação de forma contínua. 

Ademais, a iniciativa visa possibilitar o conhecimento da prevalência dos eventos adversos, o entendimento aprofundado da implicação dos efeitos dos eventos adversos nos desfechos clínicos, a mensuração da eficácia do sistema computacional no processo pedagógico e a geração de material científico de excelência para publicação. O projeto busca ativamente fomentar a produção de pesquisas acadêmicas interdisciplinares na área de segurança do paciente e garantir a inclusão prática da qualidade e segurança do cuidado nos cursos de graduação da universidade. 

Espera-se, de forma complementar, que por meio desta pesquisa e do desenvolvimento tecnológico associado sejam gerados registros de patentes de software e processos junto ao Instituto Nacional da Propriedade Industrial (INPI), artigos científicos em periódicos indexados e capítulos de livros especializados, consolidando uma produção científica e tecnológica de expressivo impacto institucional e social.

\section{Problema de Pesquisa}
\label{sec:problema}

A despeito da robustez e da validação científica internacional da metodologia IHI-GTT, identifica-se uma lacuna tecnológica substancial no que tange a plataformas computacionais integradas, abertas e customizadas ao ambiente pedagógico universitário. Os sistemas hospitalares convencionais concentram-se no faturamento ou no registro clínico documental fechado, inexistindo soluções educacionais que articulem:
\begin{enumerate}
    \item O rastreamento direcionado pelos 53 gatilhos clínicos oficiais do IHI-GTT distribuídos em 6 módulos especializados;
    \item A classificação padronizada de gravidade conforme as 9 categorias do \textit{National Coordinating Council for Medication Error Reporting and Prevention} (NCC MERP);
    \item O tratamento e análise de causa-raiz mediante ferramentas consagradas de engenharia e gestão da qualidade hospitalar (Diagrama de Ishikawa, Ciclo PDCA, Matriz GUT, Matriz SWOT, Método 5W2H e Brainstorming);
    \item Um ambiente pedagógico de simulação de prontuários clínicos que capacite estudantes de graduação de forma segura e contextualizada, sem riscos a pacientes reais.
\end{enumerate}

Nesse contexto, formula-se a seguinte questão norteadora de pesquisa:

\begin{quote}
\textit{Como conceber, arquitetar e homologar uma plataforma computacional interdisciplinar, integrando a metodologia IHI-GTT, ferramentas da qualidade hospitalar e módulos de prontuários simulados, capaz de subsidiar com alta performance, usabilidade ergonômica e rigor de engenharia de software o monitoramento de eventos adversos e o ensino clínico na Universidade Federal de Sergipe?}
\end{quote}

\section{Natureza do Projeto Departamental e Linha de Tempo}
\label{sec:natureza-projeto}

O presente documento configura-se formalmente como uma proposta de \textbf{Projeto Departamental} no âmbito do Departamento de Computação (DCOMP/CCET) da Universidade Federal de Sergipe. Desenvolvido de maneira autônoma e integral pelo discente Matheus Araujo Pereira, sob a orientação técnica do Prof. Dr. Gilton José Ferreira da Silva (DCOMP) e coorientação clínica da Profª. Drª. Ana Waleska Silva Patrício (Departamento de Enfermagem), o projeto tem sua execução planejada para o horizonte de \textbf{12 meses}.

Em conformidade estrita com as diretrizes de acompanhamento acadêmico do DCOMP para projetos departamentais, o documento estrutura-se com foco tripartite em:
\begin{itemize}
    \item \textbf{Objetivos}: formulação clara das metas gerais e específicas a serem atingidas pelo software e pela pesquisa;
    \item \textbf{Metodologia}: fundamentação teórica do IHI-GTT, taxonomia clínica, ferramentas da qualidade, arquitetura do sistema e critérios inegociáveis de engenharia de software;
    \item \textbf{Plano de Trabalho}: detalhamento cronológico de todas as fases, componentes e telas implementadas ao longo de 12 meses.
\end{itemize}

O ciclo formal de acompanhamento do projeto é balizado por dois relatórios técnicos institucionais:
\begin{enumerate}
    \item \textbf{Relatório Técnico Parcial (6 Meses)}: protocolado ao final do primeiro semestre de execução, contemplando a validação da arquitetura de segurança, o modelo de banco de dados, o catálogo completo IHI-GTT e as ferramentas de gestão da qualidade;
    \item \textbf{Relatório Técnico Final (12 Meses)}: protocolado ao término do projeto, sintetizando os resultados do módulo educacional, dashboards epidemiológicos, os dados coletados na capacitação com acadêmicos de enfermagem e os depósitos de propriedade intelectual gerados.
\end{enumerate}

O desenvolvimento do projeto é integralmente realizado por Matheus Araujo Pereira utilizando o ambiente do \textit{Antigravity IDE} sobre o sistema operacional Fedora 44 Workstation (x86\_64), aplicando agentes autônomos de inteligência artificial para maximizar a precisão arquitetural e a cobertura de testes.

\section{Motivação e Justificativa}
\label{sec:motivacao}

A motivação para o desenvolvimento deste projeto departamental decorre de três pilares fundamentais:
\begin{itemize}
    \item \textbf{Impacto Social e Assistencial}: Auxiliar a disseminação da cultura de segurança do paciente nas instituições de saúde, capacitando novos profissionais a identificarem eventos adversos com rapidez e rigor científico, migrando de uma postura meramente reativa para uma vigilância clínica proativa;
    \item \textbf{Inovação Pedagógica e Interdisciplinar com Simulação Fictícia}: Estabelecer uma cooperação científica duradoura entre o DCOMP e a Enfermagem da UFS. Para preservar total conformidade ética, o sistema opera de forma independente, sem nenhuma conexão direta com os bancos de dados reais do Hospital Universitário da UFS, baseando-se em prontuários e dados clínicos \textbf{100\% simulados e fictícios} concebidos para o treinamento dos discentes;
    \item \textbf{Excelência em Engenharia de Software}: Desenvolver um sistema aderente às mais modernas diretrizes da computação, empregando exclusivamente versões LTS de tecnologias abertas (Java 21 LTS, Spring Boot 3.3, PostgreSQL 16 LTS, Angular 18+ com Signals e TailwindCSS), submetido a um \textit{Quality Gate} estrito de 100\% de cobertura em testes e documentação.
\end{itemize}

O projeto justifica-se, portanto, pela entrega de uma solução de software robusta, gratuita, institucional e alinhada ao Programa Nacional de Segurança do Paciente (PNSP), vocacionada ao fortalecimento do ensino de graduação em Sergipe.
